\documentclass[10pt,a4paper]{amsart}
\usepackage{amsmath,amssymb,amsthm,graphicx}
\usepackage{verbatim}
\newenvironment{reference}{\comment}{\endcomment}
%\newenvironment{reference}{}{}
\newenvironment{slogan}{\comment}{\endcomment}
\newenvironment{history}{\comment}{\endcomment}

% For commutative diagrams we use Xy-pic
\usepackage[all]{xy}

% We use 2cell for 2-commutative diagrams.
\xyoption{2cell}
% Theorem environments.
%
\theoremstyle{plain}
\newtheorem{theorem}[subsection]{Theorem}
\newtheorem{proposition}[subsection]{Proposition}
\newtheorem{lemma}[subsection]{Lemma}

\theoremstyle{definition}
\newtheorem{definition}[subsection]{Definition}
\newtheorem{example}[subsection]{Example}
\newtheorem{exercise}[subsection]{Exercise}
\newtheorem{situation}[subsection]{Situation}

\theoremstyle{remark}
\newtheorem{remark}[subsection]{Remark}
\newtheorem{remarks}[subsection]{Remarks}

\numberwithin{equation}{subsection}

% Macros
%
\def\lim{\mathop{\mathrm{lim}}\nolimits}
\def\colim{\mathop{\mathrm{colim}}\nolimits}
\def\Spec{\mathop{\mathrm{Spec}}}
\def\Hom{\mathop{\mathrm{Hom}}\nolimits}
\def\Ext{\mathop{\mathrm{Ext}}\nolimits}
\def\SheafHom{\mathop{\mathcal{H}\!\mathit{om}}\nolimits}
\def\SheafExt{\mathop{\mathcal{E}\!\mathit{xt}}\nolimits}
\def\Sch{\mathit{Sch}}
\def\Mor{\mathop{\mathrm{Mor}}\nolimits}
\def\Ob{\mathop{\mathrm{Ob}}\nolimits}
\def\Sh{\mathop{\mathit{Sh}}\nolimits}
\def\NL{\mathop{N\!L}\nolimits}
\def\CH{\mathop{\mathrm{CH}}\nolimits}
\def\proetale{{pro\text{-}\acute{e}tale}}
\def\etale{{\acute{e}tale}}
\def\QCoh{\mathit{QCoh}}
\def\Ker{\mathop{\mathrm{Ker}}}
\def\Im{\mathop{\mathrm{Im}}}
\def\Coker{\mathop{\mathrm{Coker}}}
\def\Coim{\mathop{\mathrm{Coim}}}

% Boxtimes
%
\DeclareMathSymbol{\boxtimes}{\mathbin}{AMSa}{"02}

%
% Macros for moduli stacks/spaces
%
\def\QCohstack{\mathcal{QC}\!\mathit{oh}}
\def\Cohstack{\mathcal{C}\!\mathit{oh}}
\def\Spacesstack{\mathcal{S}\!\mathit{paces}}
\def\Quotfunctor{\mathrm{Quot}}
\def\Hilbfunctor{\mathrm{Hilb}}
\def\Curvesstack{\mathcal{C}\!\mathit{urves}}
\def\Polarizedstack{\mathcal{P}\!\mathit{olarized}}
\def\Complexesstack{\mathcal{C}\!\mathit{omplexes}}
% \Pic is the operator that assigns to X its picard group, usage \Pic(X)
% \Picardstack_{X/B} denotes the Picard stack of X over B
% \Picardfunctor_{X/B} denotes the Picard functor of X over B
\def\Pic{\mathop{\mathrm{Pic}}\nolimits}
\def\Picardstack{\mathcal{P}\!\mathit{ic}}
\def\Picardfunctor{\mathrm{Pic}}
\def\Deformationcategory{\mathcal{D}\!\mathit{ef}}
\usepackage{cleveref}
\crefname{originalreference}{original reference}{original references}
\Crefname{originalreference}{Original reference}{Original references}
\makeatletter
\newlabel{lemma-exact-sequence-NL}{{00S2}{source}{}{sourcepointer-1}{}}
\newlabel{lemma-exact-sequence-NL@cref}{{[originalreference][0][]00S2}{[1][0][]source}{}{}{}}
\makeatother
\crefname{originalreference}{original reference}{original references}
\Crefname{originalreference}{Original reference}{Original references}
\makeatletter
\newlabel{lemma-rig-smooth-rig-flat}{{0GH1}{source}{}{sourcepointer-2}{}}
\newlabel{lemma-rig-smooth-rig-flat@cref}{{[originalreference][0][]0GH1}{[1][0][]source}{}{}{}}
\makeatother
\crefname{originallemma}{original lemma}{original lemmas}
\Crefname{originallemma}{Original lemma}{Original lemmas}
\makeatletter
\newlabel{more-algebra-lemma-conormal-sequence-H1-regular-ideal}{{07CX}{source}{}{sourcepointer-3}{}}
\newlabel{more-algebra-lemma-conormal-sequence-H1-regular-ideal@cref}{{[originallemma][0][]07CX}{[1][0][]source}{}{}{}}
\makeatother
\crefname{originallemma}{original lemma}{original lemmas}
\Crefname{originallemma}{Original lemma}{Original lemmas}
\makeatletter
\newlabel{more-algebra-lemma-transitive-lci-at-end}{{07D4}{source}{}{sourcepointer-4}{}}
\newlabel{more-algebra-lemma-transitive-lci-at-end@cref}{{[originallemma][0][]07D4}{[1][0][]source}{}{}{}}
\makeatother
\crefname{originalsection}{original section}{original sections}
\Crefname{originalsection}{Original section}{Original sections}
\makeatletter
\newlabel{more-algebra-section-ideals}{{07CU}{source}{}{sourcepointer-5}{}}
\newlabel{more-algebra-section-ideals@cref}{{[originalsection][0][]07CU}{[1][0][]source}{}{}{}}
\makeatother
\title{Rig-smoothness kills the H\textasciicircum{}\{-1\} connecting term of the naive cotangent sequence}
\author{The Stacks Project Authors}
\date{stacks-project:restricted.tex:local-snapshot-20261012; GNU FDL 1.2 or later}
\begin{document}
\maketitle
\noindent\textit{Editorial note.} All delivered mathematical text is the author's own, copied from the pinned Stacks Project chapter file restricted.tex; no statement or proof step was rewritten, reordered or removed.\par\medskip
\noindent\textit{Editorial note.} The counted proof is the author's complete argument: the proof reduces to injectivity at primes avoiding an ideal of definition, translates it through the presentations into the intersection identity J\_q cap K\textasciicircum{}2\_q = (KJ)\_q, and derives it from rig-smooth implies rig-flat plus the H\_1-regular conormal sequence lemma\par\medskip
\noindent\textit{Editorial note.} The label lemma-exact-sequence-NL-rig-smooth names the counted statement, and the proof environment immediately following it in the source is delivered verbatim as its proof.\par\medskip
\noindent\textit{Editorial note.} The counted statement is carried by the author's own lemma environment in the chapter source.\par\medskip
\noindent\textit{Allowed standard prerequisites (editorial).}\begin{itemize}
\item the exact sequence of completed naive cotangent complexes for A -> B -> C and the induced map on H\textasciicircum{}\{-1\} (lemma-exact-sequence-NL of the same chapter, including its presentation diagram)
\item presentations of adic topologically finite type algebras B = A\{x\_1,...,x\_r\}/J, C = B\{y\_1,...,y\_s\}/J' and C = A\{x,y\}/K over the completed polynomial rings
\item rig-smooth implies rig-flat, so the ideal (K/J)\_q is generated by a regular sequence at the relevant primes (lemma-rig-smooth-rig-flat of the same chapter)
\item the conormal sequence of an ideal generated by a regular sequence is H\_1-regular (More on Algebra, Lemma more-algebra-lemma-conormal-sequence-H1-regular-ideal and its section on ideals)
\end{itemize}\medskip
\ifdefined\hypertarget\hypertarget{sourcepointer-1}{}\fi
\noindent\textit{Source reference.} Original reference 00S2: sources/stacks-restricted/pinned-original-source.tex, source lines 725--725. Prerequisite pointer invoked at proof line(s) 4568; it names a standard fact of the same chapter that the counted argument uses, and every step of the counted proof built on it is retained verbatim in the delivered body.\par\medskip
\ifdefined\hypertarget\hypertarget{sourcepointer-2}{}\fi
\noindent\textit{Source reference.} Original reference 0GH1: sources/stacks-restricted/pinned-original-source.tex, source lines 4460--4460. Prerequisite pointer invoked at proof line(s) 4594; it names a standard fact of the same chapter that the counted argument uses, and every step of the counted proof built on it is retained verbatim in the delivered body.\par\medskip
\ifdefined\hypertarget\hypertarget{sourcepointer-3}{}\fi
\noindent\textit{Source reference.} Original lemma 07CX: sources/stacks-chow/tags.txt, source lines 9392--9392. Prerequisite pointer invoked at proof line(s) 4599; it names a standard fact of a sibling chapter that the counted argument uses, and every step of the counted proof built on it is retained verbatim in the delivered body.\par\medskip
\ifdefined\hypertarget\hypertarget{sourcepointer-4}{}\fi
\noindent\textit{Source reference.} Original lemma 07D4: sources/stacks-chow/tags.txt, source lines 9399--9399. Prerequisite pointer invoked at proof line(s) 4589; it names a standard fact of a sibling chapter that the counted argument uses, and every step of the counted proof built on it is retained verbatim in the delivered body.\par\medskip
\ifdefined\hypertarget\hypertarget{sourcepointer-5}{}\fi
\noindent\textit{Source reference.} Original section 07CU: sources/stacks-chow/tags.txt, source lines 9389--9389. Prerequisite pointer invoked at proof line(s) 4602; it names a standard fact of a sibling chapter that the counted argument uses, and every step of the counted proof built on it is retained verbatim in the delivered body.\par\medskip
\begin{lemma}
\label{lemma-exact-sequence-NL-rig-smooth}
Let $A \to B \to C$ be arrows in $\textit{WAdm}^{Noeth}$
which are adic and topologically of finite type. If $B \to C$
is rig-smooth, then the kernel of the map
$$
H^{-1}(\NL_{B/A}^\wedge \otimes_B C) \to H^{-1}(\NL_{C/A}^\wedge)
$$
(see Lemma \ref{lemma-exact-sequence-NL})
is annihilated by an ideal of definition.
\end{lemma}

\begin{proof}
Let $\overline{\mathfrak q} \subset C$ be a prime ideal which does not contain
an ideal of definition. Since the modules in question are finite
it suffices to show that
$$
H^{-1}(\NL_{B/A}^\wedge \otimes_B C)_{\overline{\mathfrak q}} \to
H^{-1}(\NL_{C/A}^\wedge)_{\overline{\mathfrak q}}
$$
is injective. As in the proof of Lemma \ref{lemma-exact-sequence-NL}
choose presentations $B = A\{x_1, \ldots, x_r\}/J$,
$C = B\{y_1, \ldots, y_s\}/J'$, and
$C = A\{x_1, \ldots, x_r, y_1, \ldots, y_s\}/K$.
Looking at the diagram in the proof of Lemma \ref{lemma-exact-sequence-NL}
we see that it suffices to show that $J/J^2 \otimes_B C \to K/K^2$
is injective after localization at the prime ideal
$\mathfrak q \subset A\{x_1, \ldots, x_r, y_1, \ldots, y_s\}$
corresponding to $\overline{\mathfrak q}$. Please compare with
More on Algebra, Lemma \ref{more-algebra-lemma-transitive-lci-at-end}
and its proof. This is the same as 
asking $J/KJ \to K/K^2$ to be injective after localization
at $\mathfrak q$. Equivalently, we have to show that
$J_\mathfrak q \cap K^2_\mathfrak q = (KJ)_\mathfrak q$.
By Lemma \ref{lemma-rig-smooth-rig-flat}
we know that $(K/J)_\mathfrak q = J'_\mathfrak q$
is generated by a regular sequence.
Hence the desired intersection property follows from
More on Algebra, Lemma
\ref{more-algebra-lemma-conormal-sequence-H1-regular-ideal}
(and the fact that an ideal generated by a regular sequence
is $H_1$-regular, see
More on Algebra, Section \ref{more-algebra-section-ideals}).
\end{proof}

\end{document}
